
\section{Codes}

	\subsection{Coding with Minted*}
	
		The \inlinecode{minted} package offers customized features for adding codes. In addition, the template is designed to work seamlessly with Fira Code, a monospaced font specially crafted for programming environments.

        If you are using a desktop app as \href{https://www.texstudio.org/}{TeXstudio}, try these steps to make this package work on Windows:
		
		\begin{enumerate}
			\item Install Python — A stable version (e.g., v.3.11)
			\item Open the terminal and type \inlinecode{pip install Pygments}.
			\item Update if a newer version is available.
			\item Go to TeXStudio settings and change the default compiler to \inlinecode{pdflatex --shell-escape -interaction=nonstopmode \%.tex}.
			\item Compile and wait for the result.
		\end{enumerate}
        
		\vspace*{-9pt}
        
		\begin{tauenv}[frametitle=Caution]
			Ensure that \textbf{pygments} is properly installed and added to your system’s PATH. Otherwise, you may encounter compilation errors. Additionally, enable shell escape when compiling, as it is required for minted to process and highlight code.
		\end{tauenv}
		
		You can customize its design changing \inlinecode{\usemintedstyle} command. The different styles offered by the \href{https://ctan.org/pkg/minted}{minted package} can be preview through this link — \url{https://pygments.org/styles/}.

    \subsection{Coding with Listings}
	
		Since \inlinecode{minted} requires additional installations and can be complex in some desktop \TeX\ editors, you can use the \inlinecode{listings} package instead, which provides a simpler way to include code.
		
		\begin{code}
			\lstinputlisting[language=python]{example.py}
			\label{lst:example2}
            \caption{Python code example with listings.}
		\end{code}
		
		While \inlinecode{listings} is a simpler and widely supported package for code formatting, \inlinecode{minted} offers a more modern and powerful approach. It leverages the Pygments syntax highlighter to deliver superior coloring, language support, and styling options. One of its key advantages is the ability to easily switch between built-in color styles.
		
		In contrast, \inlinecode{listings} requires manual setup for colors, fonts, and formatting rules.
		For users who prefer fine-tuned customization, the styling options for \inlinecode{listings} are organized in \inlinecode{tau.cls} file and can be modified at any time to suit individual preferences.

	\subsection{Inline Code*}
	
		As shown in this template, inline code has a custom style for improved visual appeal. To use it, simply add \inlinecode{\inlinecode{}} command and place your code inside the braces.

